%$Id: config_nux.tex 145 2005-03-25 08:26:35Z myk $

\clearpage
\pagebreak

\bghdr{images/fond-nux}

%\begin{center}
%\includegraphics{images/logo_Linux}
%\end{center}

\subsection{Configuration sous Linux}

\subsubsection{Configuration IP}
Quelle que soit ta distribution, il faut que tu d\'ecides d'un nom de
machine --- ton pseudo en g\'en\'eral. Tu as \'egalement besoin de
conna\^itre les informations suivantes :
\begin{description}
  \item [ton adresse IP] not\'ee dor\'enavant \server{129.104.AAA.BBB} ;
  \item [l'IP de ta passerelle (\emph{gateway})] not\'ee \server{129.104.GGG.CCC} ;
  \item [ton masque de sous-r\'eseau] \server{255.255.FFF.DDD} ;
  \item [ton broadcast] qui est de la forme \server{129.104.GGG.EEE} avec $EEE = 255 - DDD$. \\
        \textcolor{red}{(attention, $GGG = AAA$ n'est valable que dans les nouveaux caserts, le PEM et le BEM. Pour Foch, Fayolle et Maunoury ce n'est donc pas n\'ecessairement vrai que $GGG =
        AAA$.)}.
\end{description}

Toutes les informations n\'ecessaires se trouvent page
\pageref{calcul_ip}.

Bien s\^ur, l'ensemble des manipulations doit se faire en tant que \emph{root}.


\paragraph{\distrib{Sous Gentoo}}

Si tu n'as jamais fait de configuration r\'eseau sur ta Gentoo, il
faudra certainement cr\'eer les fichiers qui suivent.

Le fichier \file{/etc/hostname} contient ton nom de machine. Tu peux
\'editer \file{/etc/hostname} avec ton \'editeur pr\'ef\'er\'e (\app{vi} ou
\app{vim}, \app{pico}, voire \app{emacs} si tu aimes), mais la
commande ci-dessous est suffisante :

\cmdline{echo 'ton\_pseudo' > /etc/hostname}

Le fichier \file{/etc/resolv.conf} d\'ecrit comment r\'esoudre les noms DNS. La premi\`ere ligne donne le domaine sur lequel ta machine est, ensuite
viennent les suffixes \`a utiliser par d\'efaut ; les lignes suivantes indiquent les serveurs de noms, ceux qui associent le nom \server{frankiz} ou
\server{ton\_pseudo} aux IP \server{129.104.201.51} ou \server{'ton\_ip'}. Le fichier contient donc :

\noindent \cmdline{
domain eleves.polytechnique.fr\\
search eleves.polytechnique.fr polytechnique.fr\\
nameserver 129.104.201.53\\
namaserver 129.104.201.52}

Enfin, le fichier \file{/etc/conf.d/net} contient la description de l'interface r\'eseau et les informations de routage. Il indique ton IP, ton
sous-r\'eseau, ton masque de sous-r\'eseau et la passerelle pour sortir de ton sous-r\'eseau . Bien s\^ur, remplace \texttt{eth0} par le nom de l'interface
r\'eseau que tu utilises, si ce n'est pas la carte \emph{ethernet} par d\'efaut.

\noindent \cmdline{iface\_eth0="129.104.AAA.BBB broadcast 129.104.GGG.EEE $\backslash$\\
                            netmask 255.255.FFF.DDD"\\
                            gateway="eth0/129.104.GGG.CCC"}

Tu fais le grand test en rechargeant ta configuration r\'eseau par :

\cmdline{/etc/init.d/net.eth0 restart}

puis en \emph{pingant} \fkz\ par exemple. Tu dois obtenir quelque chose comme :

\cmdline{root:~\# ping frankiz\\
PING frankiz (129.104.201.51) 56(84) bytes of data.\\
64 bytes from frankiz (129.104.201.51): ...}

\label{gentoo_mirror} Pour pouvoir utiliser \app{emerge} \`a travers le \emph{proxy} de l'\'ecole, il faut d\'efinir les variables d'environnement
ci-dessous dans le fichier \file{/etc/make.conf} :

\cmdline{http\_proxy=http://kuzh.polytechnique.fr:8080\\
ftp\_proxy=http://kuzh.polytechnique.fr:8080\\
no\_proxy=.eleves.polytechnique.fr\\
GENTOO\_MIRRORS="ftp://miroir.eleves.polytechnique.fr/gentoo http://gentoo.osuosl.org/"\\
SYNC="rsync://rsync/gentoo-portage"}

Tu peux \'evidemment ajouter d'autres miroirs (s\'epar\'es par des espaces) dans ta liste mais
\urllink{ftp://miroir} \'etant interne, il sera toujours beaucoup plus rapide que les autres. On te
conseille ici \urllink{http://gentoo.osuosl.org/} qui est un miroir tr\`es complet mais tr\`es lent.

\paragraph{\distrib{Ubuntu/Kubuntu}}
\label{Ubuntu:IP}

Lors de l'installation d'une nouvelle Ubuntu, le programme d'installation t'a normalement demand\'e de rentrer toutes ces informations de mani\`ere interactive. Si tu n'as pas fait �ce moment l\`a, tu peux les modifier comme ci-dessous avec ton \'editeur de texte pr\'ef\'er\'e (le tout avec les droits administrateurs \'evidemment!).

\begin{itemize}
\item Le fichier \emph{/etc/hostname} contient ton nom de machine. Il doit contenir uniquement:

\cmdline{tonPseudo.eleves.polytechnique.fr}

\item Le fichier \emph{/etc/resolv.conf} d\'ecrit comment associer le nom d'une machine \`a une adresse IP.
Il doit contenir:

\cmdline{
domain eleves.polytechnique.fr\\
search eleves.polytechnique.fr polytechnique.fr\\
nameserver 129.104.201.53\\
nameserver 129.104.201.52\\
}

\item Le fichier \emph{/etc/network/interfaces} contient entre autres ton IP,
ton sous-r\'eseau et la passerelle pour en sortir.
Ce fichier doit ressembler (avec \'eventuellement une config wifi \`a la suite...,
voir la page \ref{Wifi:linux}) \`a : 

\cmdline{
\# The loopback network interface\\
auto lo\\
iface lo inet loopback\\
\\
\# This is a list of hotpluggable network interfaces.\\
\# They will be activated automatically by the hotplug subsystem.\\
mapping hotplug\\
        script grep\\
        map eth0\\
\\
\# The primary network interface\\
iface eth0 inet static\\
        address   129.104.AAA.BBB\\
        netmask   255.255.255.DDD\\
        network   129.104.AAA.DDD\\
        broadcast 129.104.AAA.EEE\\
        gateway   129.104.AAA.CCC\\
}

\end{itemize}

Ensuite il faut redemarrer ta configuration r\'eseau

\cmdline{\$ sudo /etc/init.d/networking restart}

Voila ta configuration r\'eseau est termin\'ee! Tu peux la tester en pinguant frankiz

\cmdline{\$ ping frankiz}

Tu devrais voir:

\cmdline{
PING frankiz.eleves.polytechnique.fr (129.104.201.51) 56(84) bytes of data.\\
64 bytes from Frankiz.eleves.polytechnique.fr (129.104.201.51) : icmp\_seq=1 ttl=62 time=0.570 ms\\
}

{\bf Configuration du gestionnaire de paquets :}
Il faut d\'esormais configurer le gestionnaire de paquets pour qu'il utilise les miroirs du BR
et non les miroirs \`a l'ext\'erieur du campus qui sont plus lents.

Le fichier \emph{/etc/apt/sources.list} liste les miroirs utilis\'es par le gestionnaire de paquets.
Il faut commenter la premi\'ere ligne (qui correspond au cd d'installation)
ainsi que toutes les lignes non comment\'ees du fichier
(qui correspondent aux miroirs ext\'erieurs au campus) de la fa�n suivante:

\cmdline{deb cdrom:[...]/ breezy main restricted}

devient

\cmdline{\#deb cdrom:[...]/ breezy main restricted}

Il faut ensuite ajouter les lignes suivantes, qui correspondent aux miroirs du BR, au \textbf{d\'ebut} du fichier:

\cmdline{
deb ftp://miroir/ubuntu [version] main restricted universe multiverse\\
deb ftp://miroir/ubuntu [version]-backports main restricted universe multiverse\\
deb ftp://miroir/ubuntu [version]-updates main restricted universe multiverse\\
deb ftp://miroir/ubuntu [version]-security main restricted universe multiverse\\
}

o\`u \textbf{[version]} correspond \`a la version d'ubuntu install\`ee. La version actuelle est \textbf{dapper} et la pr\'ec\'edente est \textbf{breezy}.

On finit par v\'erifier que tout fonctionne en mettant \`a jour la liste des paquets disponibles:

\cmdline{\$ sudo apt-get update}

S'il n'y a pas de message d'erreur c'est que tout fonctionne nickel.


\subsubsection{Configuration antivirus \footnotesize{(elle est dr\^ole celle-l\`a hein ?)}
}

\subsubsection{Configuration du pare-feu}

La solution la plus simple pour se faire un \emph{firewall} sous linux est d'utiliser les \emph{iptables}. Pour ceci la premi\`ere \'etape est
d'installer le paquet \app{iptables} pour ta distribution. Pour savoir comment configurer ton \emph{firewall} pour le r\'eseau de l'X, consulte le Wikix.

\subsubsection{Configuration navigateur web}

\flimage{images/nux_firefox_icon}{0.12}{l} Le BR te conseille d'utiliser \app{Firefox} ou
\app{Konqueror} (le navigateur fourni par d\'efaut avec KDE). Dans tous les cas, la seule
configuration \`a effectuer est celle du \emph{proxy}. Pour \app{Firefox}, il suffit d'aller dans
\menu{Edit}, \menu{Preferences} et dans l'onglet \menu{Avanc\'e}, puis de cliquer sur l'onglet
\menu{R\'eseau}, et enfin \menu{Param\`etres de connexion} ; ensuite coche la case \menu{Adresse de
configuration automatique du proxy}, et rentre \urllink{http://frankiz/proxy.pac/}.

Sous \app{Konqueror}, cela se trouve dans le menu \menu{Configuration}, \menu{Configurer Konqueror},
dans l'onglet \menu{Serveur mandataire}.

\imagepos{images/nux_proxy_firefox}{0.55}{Configuration du proxy sous Firefox}{ht}

% \vfill
\pagebreak

\subsubsection{Configuration mail}

\flimage{images/nux_kmail_icon}{0.12}{l} Le client mail le plus
utilis\'e est \app{Kmail}, mais il en existe bien s\^ur d'autres comme
\app{Thunderbird}. La configuration est semblable, quel que soit le
client que tu utilises.

Va dans \menu{Configuration}, \menu{Configurer Kmail}. Choisis la
rubrique \menu{Comptes}. Commence par cr\'eer un nouveau compte dans
l'onglet \menu{R\'eception des messages} en cliquant sur le bouton
\menu{Ajouter\ldots} et choisis le type POP3.

\imagepos{images/nux_config_kmail_pop}{0.45}{Configuration de la r\'eception des messages sous Kmail}{pht}

Utilise les param\`etres suivants pour configurer l'onglet
\menu{G\'en\'eral} :
\begin{description}
  \item[Nom] le nom du compte, par exemple : Mails Poly
  \item[Utilisateur] rentre le login \server{poly} que t'a fourni la DSI \`a ton arriv\'ee sur le plateau
  \item[Mot de passe] et l\`a le mot de passe \server{poly}
  \item[Serveur] \server{poly.polytechnique.fr}
  \item[Port] 995
\end{description}

Ensuite, va dans l'onglet \menu{Extras} et coche la case
\menu{Utiliser SSL pour s\'ecuriser les t\'el\'echargements}.

\imagepos{images/nux_config_kmail_smtp}{0.45}{Configuration de l'envoi des messages sous Kmail}{pht}

Maintenant, dans l'onglet \menu{Envoi des messages} clique sur le
bouton \menu{Ajouter\ldots}. Utilise les param\`etres suivants pour le
configurer :
\begin{description}
  \item[Nom] le m\^eme nom de compte que pr\'ec\'edemment
  \item[Serveur] \server{poly.polytechnique.fr}
  \item[Port] 25
\end{description}
Sinon, laisse toutes les cases d\'ecoch\'ees.

\par Tu peux aussi configurer l'acc\`es \`a \app{l'annuaire LDAP} de l'\'ecole, sorte de carnet d'adresses en ligne qui contient les adresses mail de tout le monde sur le campus. Pour ce faire, commence par ouvrir \menu{Outils}, \menu{Carnet d'adresses}, puis va dans \menu{Configuration}, \menu{Configurer kAdressBook}, \menu{Consultation LDAP}. Clique ensuite sur \menu{Ajouter un h\^ote}, et configure comme suit:
\begin{description}
  \item[H\^ote] \server{ldap.eleves.polytechnique.fr}
  \item[Port] 389
  \item[Version de LDAP] 3
  \item[DN] dc=ldap,dc=eleves,dc=polytechnique,dc=fr
  \item[S\'ecurit\'e] Non
  \item[Identification] Anonyme
\end{description}
Une fois revenu dans \menu{Configuration LDAP}, coche la case \server{ldap.eleves.polytechnique.fr}. Tu as maintenant acc\`es au LDAP lors de la
r\'edaction de messages, avec tout au plus un red\'emarrage de \app{Kmail}. \imagepos{images/nux_config_ldap}{0.55}{Configuration de l'annuaire LDAP sous
Kmail}{pht}

\subsubsection{Configuration news}
\flimage{images/nux_knode_icon}{0.12}{l} Le client news le plus
utilis\'e est \app{Knode}. Parmi les autres clients news, citons
\app{Thunderbird} ou \app{slrn}. Ici aussi, la configuration est
presque ind\'ependante du logiciel choisi.

Sous \app{Knode}, c'est dans le menu \menu{Configuration}, puis \menu{Configurer Knode}. Va dans la rubrique \menu{Comptes, Forums de discussion} et
cr\'ee un compte en cliquant sur \menu{Ajouter\ldots}.

\imagepos{images/nux_config_knode}{0.45}{Configuration de Knode}{ht}

Remplis l'onglet \menu{Serveur} avec les informations suivantes :
\begin{description}
  \item[Nom] ce que tu veux pour d\'ecrire ce compte, par exemple 'News Frankiz'
  \item[Serveur] \server{frankiz.polytechnique.fr}
  \item[Port] 119
\end{description}
Ensuite occupe-toi de l'onglet \menu{Identit\'e} :
\begin{description}
  \item[Nom] mets ton pseudo dans ce champ
  \item[Organisation] X, \'ecole Polytechnique, comme tu le sens
  \item[Adresse \'electronique] ton adresse mail, pour que les gens puissent te r\'epondre par mail.
\end{description}

Enfin, pour que \app{Knode} puisse envoyer des mails, il faut aller
dans la rubrique \menu{Comptes}, sous-rubrique \menu{Serveur de
courrier (SMTP)}, et choisir comme serveur d'envoi de mails
\server{poly.polytechnique.fr}, port 25 --- c'est exactement la m\^eme
configuration SMTP que \app{Kmail}.

Si tu veux mettre une signature \`a la fin des messages que tu
posteras, il te suffit de la mettre dans l'onglet \menu{Identit\'e}.
Sur la plupart des clients la signature est interpr\'et\'ee comme
ext\'erieure au message et n'est en particulier pas incluse dans le
texte cit\'e lorsque tu r\'eponds \`a un message. Pour d\'efinir une
signature \`a la main, il suffit de mettre \verb*+-- +\ (c'est \`a dire
-{}-<espace>) sur une ligne, et tout ce qui suivra cette ligne
composera ta signature.

Il ne te reste plus qu'\`a t'inscrire \`a des \emph{newsgroups} (reporte-toi \`a la page \pageref{newsgroups} pour plus d'infos) et \`a poster !

Pour te connecter aux serveurs de \emph{news} de Polytechnique.org, qui ont un acc\`es s\'ecuris\'e, avec \app{Knode}, il y a une petite subtilit\'e car il
ne g\`ere pas le SSL. Il faut installer \app{stunnel} qui permet de d\'efinir une redirection SSL de port. Dans \file{/etc/stunnel.conf} ou
\file{/etc/stunnel/stunnel.conf} selon ta distribution, mets les lignes suivantes (les trois premi\`eres y sont en principe d\'ej\`a) :

\cmdline{\# location of pid file\\
pid = /etc/stunnel/stunnel.pid\\
\\
\# user to run as\\
setuid = stunnel\\
setgid = stunnel\\
\\
\# Use it for client mode\\
client = yes\\
\\
\# sample service-level configuration\\
\\
{[}nntps{]}\\
accept  = 1119\\
connect = ssl.polytechnique.org:563\\
TIMEOUTclose = 0
}

Il ne te reste plus qu'\`a lancer \app{stunnel} par :

\cmdline{/etc/init.d/stunnel start}

Et tu peux ainsi lire les news de Polytechnique.org en mettant \server{localhost} comme serveur et
\server{1119} comme port. Il faut aussi que tu coches \menu{Le serveur exige une identification} et
que tu rentres ton nom d'utilisateur \`a Polytechnique.org et ton mot de passe, que tu peux d\'efinir
sur \urllink{http://www.polytechnique.org/acces\_smtp.php}.
